\section{Validazione walk-forward}

\subsection{Step 10: Lancio del walk-forward}

Una volta che la strategia è implementata e lo script di walk-forward è pronto, eseguire il backtest:

\begin{lstlisting}
python tests/walk_forward_<nome_strategia>.py
\end{lstlisting}

oppure, in alternativa, lanciare il backtest dalla dashboard (tab \emph{Run Backtest}).

L'esecuzione produce:

\begin{itemize}
  \item File JSON con riepilogo metriche in \texttt{output/strategies/};
  \item File parquet con equity OOS aggregata in \texttt{output/strategies/};
  \item Log di esecuzione con eventuali warning.
\end{itemize}

\subsection{Step 11: Analisi dei risultati}

I risultati del walk-forward devono essere analizzati \emph{onestamente}, confrontandoli con i criteri di accettazione dichiarati nella issue originale.

Le metriche chiave da considerare sono:

\begin{itemize}
  \item \textbf{Sharpe ratio netto OOS}: il valore più importante. Una strategia individuale che non raggiunge almeno Sharpe 0.4-0.5 difficilmente porta valore all'ensemble;
  \item \textbf{CAGR netto OOS}: confrontato con il benchmark (SPY buy-and-hold);
  \item \textbf{Alpha vs benchmark}: il return non spiegato dal beta di mercato;
  \item \textbf{MaxDD}: stabilità della curva, no drawdown catastrofici;
  \item \textbf{Correlazione con stream esistenti}: se la nuova strategia ha correlazione superiore a 0.6 con uno stream già presente, il suo contributo marginale all'ensemble è limitato;
  \item \textbf{Comportamento in periodi di stress}: 2008, 2020, 2022. La strategia non deve crollare drasticamente in regimi avversi;
  \item \textbf{Stabilità nel tempo}: confrontare lo Sharpe nei primi anni del backtest con gli ultimi anni. Una strategia il cui edge si erode significativamente è un campanello d'allarme.
\end{itemize}

\subsection{Step 12: Decisione di accettazione o scarto}

A questo punto si presenta un bivio:

\begin{itemize}
  \item Se la strategia rispetta i criteri di accettazione della issue, si procede con l'\textbf{esito positivo} descritto nella Sezione \ref{sec:esito_positivo};
  \item Se la strategia non rispetta i criteri (Sharpe troppo basso, MaxDD eccessivo, instabilità inter-temporale, correlazione troppo alta con altri stream, ecc.), si procede con l'\textbf{esito negativo} descritto nella Sezione \ref{sec:esito_negativo}.
\end{itemize}

\subsubsection{Anti-pattern: salvare a tutti i costi}

In presenza di risultati insoddisfacenti, è tentante "aggiustare" la strategia con micro-ottimizzazioni ai parametri fino a ottenere numeri presentabili. Questo è curve-fitting e va evitato. Le iterazioni accettabili sono:

\begin{itemize}
  \item Correggere bug nell'implementation (es. lookahead bias, errata gestione di dividendi);
  \item Aggiungere features mancanti richieste dalla logica originale della strategia (es. filtri di liquidità);
  \item Cambiare l'asset class su cui opera la strategia se il razionale lo giustifica.
\end{itemize}

Le iterazioni non accettabili sono:

\begin{itemize}
  \item Cambiare lookback finché il Sharpe migliora;
  \item Escludere asset post-hoc perché hanno avuto performance peggiori;
  \item Aggiungere filtri di regime tunati sui dati di test;
  \item Restringere il periodo di backtest a quello in cui i numeri sembrano migliori.
\end{itemize}

Se ci si trova a iterare in queste direzioni, la strategia probabilmente non ha un edge robusto. Procedere all'esito negativo è preferibile a integrare nel sistema una strategia overfittata.
